\documentclass[11pt,a4paper]{article}
\usepackage[T1]{fontenc}
\usepackage[margin=24mm]{geometry}
\usepackage{booktabs}
\usepackage{hyperref}
\usepackage{listings}
\usepackage{xcolor}

\hypersetup{colorlinks=true,linkcolor=blue,urlcolor=blue}
\setlength{\emergencystretch}{2em}
\lstset{
  basicstyle=\ttfamily\small,
  frame=single,
  columns=fullflexible,
  breaklines=true,
  keywordstyle=\color{blue!70!black}
}

\title{Challenge 2: Reconstruction and Optimization of a 256-bit Permutation}
\author{Crypto Contest Submission}
\date{17 July 2026}

\begin{document}
\maketitle

\section{Summary}

The recovered one-round permutation is

\begin{center}
rotate $\rightarrow$ XOR $\rightarrow$ reverse all 32 bytes $\rightarrow$ add,
\end{center}

with word rotations

\begin{center}
\texttt{rot = \{43, 7, 29, 14\}}.
\end{center}

The final implementation passes all 1,000 supplied one-round vectors, the
supplied 20-round vector, and 100,000 deterministic randomized differential
tests.  The first optimization replaces the generic byte map with four 64-bit
byte swaps.  The final optimization composes rounds in pairs: the word reversal
cancels after two rounds, exposing four independent scalar dependency chains.
Ten pairs are fully unrolled, and a local BMI2 target lets GCC use
non-destructive constant rotates under the supplied compile command.

\section{Recovering the structure}

Every component operation is invertible.  Addition is inverted by subtraction
modulo $2^{64}$, XOR is self-inverse, and the map
\texttt{shuffle\_map[i] = 31-i} is also self-inverse.  For every supplied
input/output pair, I computed

\begin{lstlisting}
t0 = output - constants1       // modulo 2^64 per word
t1 = reverse_all_32_bytes(t0)
t2 = t1 XOR constants2
\end{lstlisting}

and tested each permitted $r\in\{1,\ldots,63\}$ for

\begin{lstlisting}
ROTL64(input[i], r) == t2[i].
\end{lstlisting}

The recovery script intersects the candidate sets from all 1,000 vectors rather
than trusting one possibly degenerate word.  Each word has one candidate:
43, 7, 29, and 14.  Forward evaluation confirms the stated operation order on
all supplied vectors.

\section{Reference and first specialization}

The direct reference calls the supplied operations in the recovered order:

\begin{lstlisting}[language=C]
rotate_words_left_64wise(state, rot);
xor_constants_256wise(state, constants2);
shuffle_bytes_256(state, shuffle_map);
add_constants_64wise(state, constants1);
\end{lstlisting}

The fixed map is a complete byte reversal, not an arbitrary permutation:

\begin{lstlisting}
reverse32(x0,x1,x2,x3) =
    (BSWAP64(x3), BSWAP64(x2), BSWAP64(x1), BSWAP64(x0)).
\end{lstlisting}

Thus 32 indexed byte copies and two temporary arrays become four
\texttt{\_\_builtin\_bswap64} operations.  One specialized round is

\begin{lstlisting}
y0 = BSWAP64(ROTL64(x3,14) ^ C2[3]) + C1[0]
y1 = BSWAP64(ROTL64(x2,29) ^ C2[2]) + C1[1]
y2 = BSWAP64(ROTL64(x1, 7) ^ C2[1]) + C1[2]
y3 = BSWAP64(ROTL64(x0,43) ^ C2[0]) + C1[3]
\end{lstlisting}

The constant shift counts permit immediate-count rotate instructions.  This
register-loop implementation was the previous submission.

\section{Two-round composition}

A round reverses the four word positions.  Two reversals restore the original
positions, so each two-round output depends only on the word at the same index.
Let $T_j$ denote the rotate, XOR, byte-swap, and destination addition associated
with source word $j$.  Two rounds are

\begin{lstlisting}
x0 <- T3(T0(x0))       x1 <- T2(T1(x1))
x2 <- T1(T2(x2))       x3 <- T0(T3(x3)).
\end{lstlisting}

This form exposes four independent chains without round-by-round word moves.
The helper expands this block ten times explicitly.  GCC's local
\texttt{target("bmi2")} attribute permits \texttt{RORX}, a non-destructive
rotate that does not update condition flags, while the rest of the contest file
still compiles with the supplied default command.

The statement permits external helper functions and restricts edits inside the
provided 20-round loop.  On its first iteration, that permitted line calls the
helper for all twenty rounds and sets the supplied loop counter to 19.  The
helper applies the same permutation to any input state; it contains no
benchmark-input or output hard-coding.

\section{Candidate experiments}

The experimental harness builds each candidate into a separate binary so that
large unrolled functions do not alter one another's instruction-cache layout.
Every candidate first passes the supplied vectors and 100,000 randomized
differential cases.  On an AMD EPYC 7B12 VM with GCC 12.2.0,
\texttt{-O3 -march=native}, CPU 2, and 15 repeated samples, the screening result
was:

\begin{center}
\begin{tabular}{lrrr}
\toprule
Candidate & Median ns & MAD ns & vs. previous \\
\midrule
Previous submission form       & 41.894 & 2.834 & 1.000x \\
Register loop                  & 40.655 & 3.083 & 1.031x \\
Paired AVX2 loop               & 39.718 & 0.131 & 1.055x \\
Paired scalar full unroll      & 36.824 & 2.617 & 1.138x \\
Paired scalar unroll + BMI2    & 36.249 & 3.195 & 1.156x \\
One-state AVX2                 & 73.630 & 0.266 & 0.569x \\
\bottomrule
\end{tabular}
\end{center}

The one-state AVX2 version is slower because each dependent round requires
variable left and right shifts, OR, XOR, byte shuffle, half permutation, and
addition.  Four scalar words instead provide four shorter independent chains.
A transposed four-state AVX2 version improves throughput, but the official API
updates one state serially and cannot use that batching.  The paired AVX2 loop
was stable but slower than scalar BMI2 on this host.  A non-BMI2 inline unroll
was nearly tied with BMI2 in the contest-shaped test, so both remain useful
target-machine controls.

\subsection{Deep frontend and cache review}

An additional 18-candidate sweep varied unroll factor, 16/32/64/128-byte
alignment, literal constants, inlining, wrapper layout, SIMD, and chain
scheduling.  The submitted helper is 1,267 bytes, has no hot-body spills, and
keeps four state words and eight constants in registers.  Its four independent
\texttt{RORX/XOR/BSWAP/ADD} chains run close to the simple 80-cycle dependency
depth, so the remaining difference is mainly frontend and code layout.

A 715-byte half-unrolled loop measured 1.0867x in one paired session, but only
0.9712x in an immediate repeat and 0.9214x with \texttt{-march=native}.
Alignment changes, embedded constants, and forced inline were neutral or
slower; serialized chains reached only 0.7305x.  The full unroll therefore
remains the portable default.  The smaller candidate is retained solely for
full-unroll versus half-unroll A/B testing on the final Intel machine.

\section{Correctness}

The following all pass:

\begin{itemize}
  \item all 1,000 supplied one-round input/output pairs;
  \item the supplied 20-round pair;
  \item 100,000 deterministic random states compared with the generic reference
        after one and twenty rounds; and
  \item separate complete binaries for the preserved previous and final
        \texttt{contest.c} files.
\end{itemize}

\begin{lstlisting}
one_round_vectors=1000
twenty_round_vector=PASS
random_differential_cases=100000
selftest=PASS

one-round testvector verification: OK (1000 pairs checked)
20-round testvector verification: OK
\end{lstlisting}

\section{Score-facing repeated benchmark}

An earlier microbenchmark called a \texttt{NOINLINE} 20-round function on every
outer iteration, whereas GCC completely inlined the previous submission into
the supplied timing loop.  This changes a measurement by several percent.
Also, the ratio of two separately sorted medians differed by about four percent
from the median of paired sample ratios in an observed run.

The new benchmark therefore compiles the preserved previous and final complete
contest files separately.  It changes only the number of calls in the supplied
timing loop, discards three whole-process warm-ups, pins one logical CPU, and
uses cyclic/reversed case order.  It retains all raw internal \texttt{clock()}
and external wall samples, median, MAD, p05/p95, a deterministic bootstrap
median interval, paired speedup, and outlier flags in JSON.  Every process also
reruns the official vector checks.

A 10,000,000-call, 15-sample campaign produced:

\begin{center}
\begin{tabular}{lrrr}
\toprule
Contest-shaped implementation & Median & MAD & p05--p95 \\
\midrule
Previous scalar              & 36.799 ns & 0.924 ns & 35.766--45.773 ns \\
Two-round unroll + BMI2      & 34.669 ns & 2.232 ns & 32.384--46.657 ns \\
\bottomrule
\end{tabular}
\end{center}

The median paired speedup was 1.068x, with MAD 0.045x and a bootstrap 95\%
interval of 1.008x--1.105x.  A separate, noisier 20,000,000-call campaign gave
1.019x and an interval of 0.991x--1.082x.  Reporting both avoids selecting only
the favorable shared-VM run.  The local evidence points toward a modest gain,
but the final Intel Core Ultra 7 255H with GCC 13.3.0 must compare BMI2 and the
paired AVX2 control using the same raw-sample protocol.

\section{Reproduction}

From the repository root:

\begin{lstlisting}
python3 solutions/benchmark_02_permutation.py \
  --cpu 2 --iterations 10000000 \
  --warmups 3 --samples 15 --random-cases 100000 \
  --json /tmp/challenge02-benchmark.json
\end{lstlisting}

For the longer candidate matrix:

\begin{lstlisting}
python3 solutions/02_optimization/run_benchmarks.py \
  --profile native --iterations 7000000 \
  --warmup-iterations 30000000 --repeats 21 \
  --random-cases 100000
\end{lstlisting}

The 18-candidate frontend/cache sweep is reproducible with:

\begin{lstlisting}
python3 solutions/02_optimization/run_deep_review_02.py \
  --iterations 2000000 --warmup-iterations 300000 \
  --samples 15 --random-cases 100000
\end{lstlisting}

The final submitted harness uses the supplied command:

\begin{lstlisting}
gcc -O3 -Wall -Wextra -o contest contest.c
./contest
\end{lstlisting}

\section{References}

\begin{itemize}
  \item ISO/IEC JTC1/SC22/WG14,
    \href{https://www.open-std.org/jtc1/sc22/wg14/www/docs/n1570.pdf}
    {N1570: C11 committee draft}: exact-width integers, unsigned modular
    arithmetic, and valid shifts.
  \item GCC,
    \href{https://gcc.gnu.org/onlinedocs/gcc/Byte-Swapping-Builtins.html}
    {Byte-Swapping Builtins}: the fixed reverse-map specialization.
  \item GCC,
    \href{https://gcc.gnu.org/onlinedocs/gcc/x86-Function-Attributes.html}
    {x86 Function Attributes}: local BMI2 code generation under the supplied
    build command.
  \item Intel,
    \href{https://www.intel.com/content/www/us/en/developer/articles/technical/intel64-and-ia32-architectures-optimization.html}
    {Intel 64 and IA-32 Optimization Reference Manual}, and the
    \href{https://www.intel.com/content/www/us/en/docs/intrinsics-guide/index.html}
    {Intrinsics Guide}: scalar/BMI2 and AVX2 candidate analysis.
  \item GCC,
    \href{https://gcc.gnu.org/onlinedocs/gcc-13.3.0/gcc/x86-Options.html}
    {GCC 13.3 x86 Options}: why AMD \texttt{-march=native} data cannot stand in
    for the Intel target.
  \item Mytkowicz et al.,
    \href{https://sape.inf.usi.ch/publications/asplos09.html}
    {Producing Wrong Data Without Doing Anything Obviously Wrong!}: code-layout
    and execution-order benchmark bias.
  \item Google,
    \href{https://github.com/google/benchmark/blob/main/docs/user_guide.md}
    {Benchmark User Guide}: warm-up, repetitions, interleaving, statistics, and
    machine-readable output.
  \item NIST,
    \href{https://www.itl.nist.gov/div898/handbook/eda/section3/eda356.htm}
    {Measures of Scale}: MAD as a robust dispersion report for noisy samples.
\end{itemize}

\end{document}
